\section{Discussion}
\label{sec:discussion}

\paragraph{Available and realized amplification.}
The factorization separates two questions that are often merged. The first is how much amplification the learned dynamics make available, measured by $G_{\max}$, by the forced gain $\|\mathcal L_T\|_2$ or, along a forecast, by $\bar g_T$. The second is how much of it the defects of the surrogate use. In the current experiments, the worst-case estimate $B_T$ exceeds the realized error by a factor of about two at $0.3$ Lyapunov times and by a factor of five to seven at about one Lyapunov time, and the difference is removed by the efficiency and the cancellation in roughly comparable parts. The low efficiency does not indicate that the defects avoid the amplified directions, since random defects with the same second moment realize the same fraction of the available gain, and it decreases with the horizon as the stable rank of the tail propagators (Proposition~\ref{prop:nulls}). The gain that matters for the error of a typical defect is therefore the covariance-weighted gain of generalized stability theory \citep{farrell1996a,farrell1996b} rather than $\sigma_1$, and a bound based on $\sigma_1$, such as the Lipschitz-type bounds commonly used to discuss compounding errors, is pessimistic by a factor that grows with the dimension of the unstable subspace explored by the defects.

\paragraph{Coherence of the defects.}
The transported defects add much more coherently than independent vectors, with coherence factors of four to five at about one Lyapunov time, against $1.05$ to $1.25$ for independent defects with the bias of the surrogates. This is the learned-surrogate counterpart of the quadratic short-time growth of the model error described by \citet{vannitsem2002}, extended to strongly anisotropic tail propagators. Since the coherence comes from the temporal correlation of the defects and not from their mean, removing a constant bias from a surrogate would change little. Objectives that only penalize the size of individual defects, or of individually transported defects, do not act on the cross terms $\mathbf q_j^\top\mathbf q_k$ that make the error grow faster than a random walk, whereas multistep objectives do act on them, since $\|\delta^{\rm lin}_T\|_2^2$ contains all cross terms.

\paragraph{Which factor makes a forecast worse.}
The variance attribution answers a different question from the median budget. Below a few tenths of a Lyapunov time, the defect magnitude explains most of the differences between forecasts. Beyond that, they are explained by the available gain and the coherence and, for the FNO, by the state-to-state variation of the alignment between defects and amplified directions, which is invisible in the median efficiency and appears only in the variance attribution. The worst-case gain, which is the quantity targeted by Jacobian-based regularization \citep{pan2018,nie2026,pervez2026}, becomes the largest single contribution only near and beyond one Lyapunov time, when the first-order description starts to lose accuracy.

\paragraph{Surrogate accuracy.}
Longer training reduces the defect magnitude by about one order of magnitude and extends the first-order regime, but changes the dimensionless factors only moderately. The main qualitative difference is in the nonlinear factor, since the L96-20 surrogates leave the attractor in free run and their errors at about one Lyapunov time already contain an amplitude drift. A surrogate with accurate one-step predictions and accurate Jacobians along the reference trajectory can therefore still have unstable free rollouts \citep{park2024}, and the budget identifies this case through $N_T$.

\paragraph{Multistep fine-tuning.}
The paired attribution shows that the benefit of the multistep objective for the L96-20 surrogates comes from the worst-case gain, that is, from tangent dynamics that become less expansive than those of the reference system, and not from smaller or better oriented defects. Damping the learned dynamics reduces the error of an under-trained surrogate whose free rollouts drift away from the attractor, but it moves the tangent dynamics further from the reference ones, which may degrade other quantities, such as the statistics of long free runs. Surrogates that already reproduce the reference tangent dynamics leave no such margin, and for them the multistep objective changes neither the error nor any factor of the budget. This does not imply that multistep training is ineffective in general, since longer horizons, training from scratch or larger models may behave differently \citep{brandstetter2022,list2025}. It does show that an improvement obtained with a multistep objective should be checked against a converged one-step baseline, and the factorization indicates which part of the budget the objective changed.

\paragraph{Toward prospective use.}
The ensemble estimate of Section~\ref{sec:res_ensemble} turns the retrospective account into a prospective one, and the analysis shows why scores built from the defect statistics fail. The realized error depends on the actual defects along the forecast, through their magnitude in the region of the attractor visited by the forecast, their temporal coherence and their orientation, which no statistic of the defect distribution contains. Independently trained surrogates disagree where they are inaccurate, and their disagreement reproduces these properties well enough to rank the forecasts, provided that the members make sufficiently independent errors, which fails for the under-trained L96-20 surrogates. Ensemble spreads are routinely used as error estimates \citep{leutbecher2008,lakshminarayanan2017}, and the propagated proxy $\hat L_T$ adds a first-order estimate that can be decomposed in the same way as the realized error.

\paragraph{Reference and model propagators.}
Appendix~\ref{app:stress} confirms that the leading right singular vector of $\Phim$ predicts the worst finite-amplitude perturbation of the learned rollout up to sizable amplitudes. This is compatible with the results above: the worst-case direction describes the most dangerous perturbation, but the defects of a surrogate are not drawn from it, and the realized error depends on the full singular spectrum, on the defect covariance and on the coherence of the transported defects.
